% ============================================================
% FUTURE WORK
% (Draft — to be expanded after Discussion is finalized)
% ============================================================

\chapter{Future Work}
\label{ch:future_work}

This chapter outlines promising research directions that build on the findings and address the limitations of this thesis.


% ============================================================
\section{Combining Online Calibration with Adaptive NCM}
\label{sec:fw_combined}
% ============================================================

The most direct extension of this work is combining the two independent improvements demonstrated in this thesis: online calibration management (Experiment~2) and adaptive NCM via normalized residuals (Experiment~3). This thesis deliberately kept these components separate to isolate their individual contributions, but a practical deployment would benefit from both simultaneously.

Specifically, an online conformal prediction system with a normalized NCM ($R_i = |y_i - \hat{y}_i| / \sigma_i$) where both the calibration set and the difficulty estimates ($\sigma_i$) are updated over time would:
\begin{itemize}
    \item Maintain coverage validity through fresh calibration data (as shown in Experiment~2),
    \item Improve interval efficiency through sample-specific widths (as shown in Experiment~3), and
    \item Potentially close the remaining gap between the best online method (PICP = 74.6\%) and the 90\% target.
\end{itemize}

\paragraph{Preliminary experimental evidence.} To verify this hypothesis, we conducted a preliminary experiment combining online calibration strategies with both absolute and normalized NCM functions at the segment level, aggregated to route level via simple summation. Table~\ref{tab:fw_combined_results} reports the results.

\begin{table}[htbp]
\centering
\caption{Preliminary results: combining online calibration with NCM variants (segment-level CP aggregated to route level via simple summation, 90\% target).}
\label{tab:fw_combined_results}
\begin{tabular}{lrrrr}
\toprule
\textbf{Method} & \textbf{PICP} & \textbf{MPIW (s)} & \textbf{Cal.\ Err} & \textbf{Winkler} \\
\midrule
Static + Absolute NCM              & 0.8871 & 3998.7 & 0.0129 & 5714.7 \\
Static + Normalized NCM            & 0.9241 & 3865.8 & 0.0241 & 4878.1 \\
Online Expanding + Absolute NCM    & 0.8914 & 4084.5 & 0.0086 & 5706.4 \\
Online Sliding-14d + Absolute NCM  & 0.8939 & 4122.4 & 0.0061 & 5706.9 \\
Online Expanding + Normalized NCM  & 0.9115 & 3744.2 & 0.0115 & 4857.2 \\
Online Sliding-14d + Normalized NCM & 0.9053 & 3678.5 & 0.0053 & 4848.7 \\
\bottomrule
\end{tabular}
\end{table}

Three findings emerge from this preliminary experiment. First, the normalized NCM improves interval efficiency regardless of the calibration strategy: every normalized variant achieves a lower Winkler score than its absolute counterpart, confirming that adaptive interval widths and calibration freshness are complementary rather than redundant. Second, the combination of online calibration with normalized NCM achieves the best overall profile: Online Sliding-14d with Normalized NCM reaches PICP~=~0.9053 with the lowest calibration error (0.0053) and the best Winkler score (4848.7), demonstrating that both components contribute independently to interval quality. Third, even the combined approach produces route-level MPIW values in the 3600--4100~second range ($\approx$60--68~minutes), which, while improved over the 4514~seconds of the static aggregated-sum baseline in Experiment~3, remains operationally wide. This residual width is a structural consequence of summing independent segment intervals over 40+ segments, and further reduction requires the covariance-aware aggregation and stronger base models discussed in Sections~\ref{sec:fw_covariance} and~\ref{sec:fw_deep_learning}.

The key open question remains whether a fully integrated system---online calibration, normalized NCM, and covariance-aware aggregation---can bring route-level MPIW below 10~minutes while maintaining valid coverage. This integration represents the natural next step toward a deployment-ready framework.


% ============================================================
\section{Alternative NCM Functions}
\label{sec:fw_ncm}
% ============================================================

This thesis explored only two NCM functions (absolute and normalized residuals). Several alternative NCM designs could improve interval efficiency:

\begin{itemize}
    \item \textbf{Quantile regression-based NCM}: Train the model to directly predict conditional quantiles rather than the mean, then use the quantile range as the NCM. This naturally produces asymmetric intervals that reflect the right-skewed nature of bus travel time distributions.

    \item \textbf{Conformal quantile regression (CQR)}: A hybrid approach where a quantile regression model is conformalized, combining the flexibility of quantile regression with the coverage guarantee of conformal prediction.

    \item \textbf{Locally-weighted NCM}: Use a $k$-nearest-neighbors approach to estimate local residual variability, producing NCM values that adapt to the feature space rather than only to the segment identity.

    \item \textbf{Time-dependent NCM}: Design an NCM that explicitly incorporates the temporal distance from the calibration period, down-weighting older calibration samples.

    \item \textbf{Conditional quantile thresholds}: Rather than computing a single global quantile $\hat{q}$ from the entire calibration set, compute separate quantiles for operationally distinct subgroups---for example, one $\hat{q}$ for peak hours and another for off-peak, or one for weekdays and another for weekends. The conditional coverage analysis in Experiment~1 (Table~T4.3) showed that coverage degrades unevenly across time-of-day categories, with midday and evening periods suffering the most. Condition-specific quantiles would tighten intervals during predictable periods (off-peak, weekend) and widen them only where the data warrants it. This approach has a direct theoretical grounding: conditional conformal prediction replaces the marginal exchangeability assumption with a group-conditional one, which is easier to satisfy within homogeneous subgroups. A practical implementation could stratify the calibration set by time period and day type, compute per-stratum quantiles, and apply the appropriate quantile at prediction time based on the input features. The expected benefit is a meaningful reduction in average MPIW with no loss in per-group coverage validity.
\end{itemize}

Each of these alternatives trades simplicity for potential efficiency gains, and the optimal choice likely depends on the specific transit network and data characteristics.


% ============================================================
\section{Deep Learning Baselines}
\label{sec:fw_deep_learning}
% ============================================================

This thesis uses XGBoost as the sole baseline model. Extending the framework to deep learning architectures (LSTM, GRU, Transformer, Graph Neural Networks) would provide two benefits:

\begin{enumerate}
    \item \textbf{Improved point prediction and narrower intervals}: Sequential models such as LSTMs or Temporal Fusion Transformers can capture temporal dependencies (e.g., propagating congestion, signal phase cycles) that tree-based models treat as noise. Stronger point predictions produce smaller residuals, which directly reduce the difficulty estimates ($\sigma_i$) used by the normalized NCM. Since interval width is proportional to $\hat{q} \times \sigma_i$, a model that halves the typical residual magnitude would roughly halve the interval width---potentially bringing the aggregated route-level MPIW from the current 60--75~minutes into the 15--30~minute range, which is far more operationally useful.
    \item \textbf{Different residual structure}: Deep learning models may produce differently distributed residuals (e.g., more or less heteroscedastic), which would test whether the findings about coverage degradation under drift are model-independent or model-specific.
\end{enumerate}

A particularly relevant comparison would be to test whether the conformal prediction layer ``equalizes'' different base models --- that is, whether a weak model with good calibration can achieve similar practical performance to a strong model with poor calibration. The model-agnostic nature of conformal prediction makes this comparison straightforward: the entire CP layer (calibration, online updates, segment decomposition) developed in this thesis can be applied to any regression model without modification.


% ============================================================
\section{Covariance-Aware Segment Aggregation}
\label{sec:fw_covariance}
% ============================================================

The simple summation aggregation used in Experiment~3 treats segment intervals as independent, summing lower and upper bounds directly:
\begin{equation}
    W_{\text{route}} = \sum_{i=1}^{N} W_i
\end{equation}
where $W_i = 2 \hat{q} \cdot \sigma_i$ is the interval width of segment $i$. This independence assumption is conservative: when consecutive segments share the same traffic conditions (e.g., both lie on the same congested corridor), their prediction errors are positively correlated, and the simple sum overestimates the true route-level uncertainty.

A covariance-aware aggregation would model this inter-segment correlation explicitly. Under a multivariate framework, the route-level uncertainty can be expressed as:
\begin{equation}
    W_{\text{route}}^{\text{cov}} = 2 \hat{q} \cdot \sqrt{\sum_{i=1}^{N} \sigma_i^2 + 2 \sum_{i < j} \rho_{ij} \sigma_i \sigma_j}
    \label{eq:cov_aggregation}
\end{equation}
where $\rho_{ij}$ is the correlation between the residuals of segments $i$ and $j$. When correlations are positive but less than one (as expected in practice), the square root produces a route-level width that is substantially smaller than the linear sum. To illustrate: for $N = 42$ segments with equal $\sigma$ and uniform pairwise correlation $\rho = 0.3$, the ratio of the covariance-aware width to the simple sum is:
\begin{equation}
    \frac{W^{\text{cov}}}{W^{\text{sum}}} = \frac{\sqrt{N + N(N-1)\rho}}{N} = \frac{\sqrt{42 + 42 \times 41 \times 0.3}}{42} \approx 0.56
\end{equation}
This would reduce the aggregated MPIW from approximately 4500~seconds to approximately 2500~seconds ($\approx$42~minutes)---a meaningful improvement toward operationally useful intervals.

The practical challenge lies in estimating the $N \times N$ correlation matrix $\rho_{ij}$ reliably from the calibration data. Two estimation strategies are natural candidates:
\begin{itemize}
    \item \textbf{Empirical correlation from calibration residuals}: Compute the sample correlation matrix of segment residuals across calibration trips. This is straightforward but requires sufficient calibration data per segment pair. With $N = 42$ segments, only $42 \times 41 / 2 = 861$ unique pairs need to be estimated, which is feasible given several thousand calibration trips.
    \item \textbf{Structured correlation models}: Impose a spatial structure on the correlation matrix --- for instance, assuming that correlation decays with the number of intermediate segments ($\rho_{ij} = \rho^{|i-j|}$ for some base correlation $\rho$). This reduces the number of free parameters to one and regularizes the estimate, at the cost of assuming a specific correlation structure.
\end{itemize}

Integrating covariance-aware aggregation with the online calibration and normalized NCM framework would produce a system where interval widths reflect not only per-segment difficulty but also the spatial correlation structure of the route, yielding tighter and more realistic route-level intervals without sacrificing coverage validity.


% ============================================================
\section{Multi-City Generalization}
\label{sec:fw_multi_city}
% ============================================================

The current study is limited to three routes in Astana, Kazakhstan. Extending the framework to multiple cities with different characteristics (e.g., European cities with tram networks, Asian megacities with high congestion variability, developing-world cities with less predictable infrastructure) would test the generalizability of the findings. In particular:

\begin{itemize}
    \item Do different cities exhibit different drift speeds? If so, does the optimal sliding window size change?
    \item Are the same segments consistently uncertain across different seasons and cities?
    \item Can transfer learning be used to bootstrap calibration in a new city using data from a similar city?
\end{itemize}


% ============================================================
\section{Real-Time Deployment Architecture}
\label{sec:fw_deployment}
% ============================================================

This thesis evaluates all methods in an offline, batch-processing setting. A real-time deployment would introduce additional considerations:

\begin{itemize}
    \item \textbf{Latency constraints}: The hourly CP updates require re-creation of the \texttt{WrapCalibratedExplainer} object, which takes approximately 15--20~ms per prediction. For a fleet of hundreds of buses with per-stop predictions, this may require optimization.
    \item \textbf{Streaming calibration}: Rather than batch updates (daily or hourly), a truly streaming approach could update the calibration set after each individual observation, maintaining a sliding window in real-time.
    \item \textbf{Anomaly detection}: In production, the system should detect and flag anomalous observations (e.g., GPS errors, route detours) before they contaminate the calibration set.
    \item \textbf{Model retraining triggers}: The framework could be extended with automatic triggers for model retraining when the prediction error exceeds a threshold, combining model adaptation with calibration adaptation.
\end{itemize}


% ============================================================
\section{External Variable Integration}
\label{sec:fw_external}
% ============================================================

Weather conditions, special events, school schedules, and real-time traffic data are known to affect bus travel times but were not available in the dataset used in this thesis. Integrating these variables could:

\begin{itemize}
    \item Reduce the residual variability captured by the NCM (by explaining some of the variance through features rather than treating it as noise),
    \item Enable event-aware calibration strategies (e.g., using separate calibration sets for rainy vs.\ dry days),
    \item Provide a more realistic test of how the framework handles sudden, event-driven distribution shifts as opposed to the gradual seasonal drift observed in the current dataset.
\end{itemize}


% ============================================================
\section{Online Segment-Level Decomposition}
\label{sec:fw_online_segment}
% ============================================================

Experiment~3 demonstrated segment-level decomposition with a static calibration set. Combining this with online calibration (Experiment~2) at the segment level would create a fully adaptive, spatially interpretable uncertainty framework. The per-segment difficulty estimates ($\sigma_s$) could also be updated over time, allowing the system to detect when a previously stable segment becomes a new uncertainty hotspot (e.g., due to road construction or a new traffic signal).

\paragraph{Route-level evidence.} The preliminary experiment reported in Table~\ref{tab:fw_combined_results} (Section~\ref{sec:fw_combined}) already demonstrates the benefit of combining online calibration with normalized NCM at the \emph{route} level: the Online Sliding-14d + Normalized NCM configuration achieves a calibration error of only 0.0053 with a Winkler score of 4848.7, substantially outperforming all single-component baselines. However, this route-level combination still produces intervals of 3600--4100~seconds because it operates on a single aggregated travel time rather than decomposing uncertainty across segments.

\paragraph{Expected benefits of the segment-level extension.} Applying the same online + normalized combination at the segment level and aggregating via simple summation would address three limitations simultaneously:
\begin{enumerate}
    \item \textbf{Reduced over-coverage}: Experiment~3's static segment-level CP achieves 99.21\% route-level coverage---far above the 90\% target. This over-coverage implies that the per-segment quantile thresholds are larger than necessary. Online calibration would keep these thresholds current, reducing over-coverage toward the 90\% target and thereby \emph{narrowing} the aggregated route-level intervals. The reduction from 99\% to 90\% coverage alone could shrink MPIW substantially, because the 99th-percentile quantile is considerably larger than the 90th-percentile quantile of the nonconformity score distribution.
    \item \textbf{Drift-adaptive segment-level widths}: Under static calibration, a segment whose prediction difficulty increases over the test period (e.g., due to new construction) retains its Week~4 difficulty estimate. Online updates would detect this change and widen the interval for that specific segment, maintaining local coverage validity without inflating the intervals of unaffected segments.
    \item \textbf{Full interpretability under drift}: The uncertainty attribution, feature explanations, and spatial maps demonstrated in Experiment~3 would remain available, but now with attributions that reflect \emph{current} conditions rather than the potentially stale Week~4 calibration. Transit operators would receive up-to-date information about where uncertainty concentrates and why.
\end{enumerate}

This is arguably the most impactful future direction, as it would deliver a complete, deployment-ready system that maintains valid coverage, provides efficient intervals, and explains where uncertainty originates---all while adapting to changing conditions in real time. Combined with the covariance-aware aggregation described in Section~\ref{sec:fw_covariance} and a stronger base model (Section~\ref{sec:fw_deep_learning}), this integrated system has the potential to bring route-level MPIW into the 5--15~minute range while preserving both coverage validity and spatial interpretability.